\chapter{Testing}\label{ch:testing}

\index{testing}\index{clat test}\index{test blocks}

Untested code is code that's waiting to betray you.
You might trust it today---it ran once, it printed the right number---but refactor a function, add a parameter, swap a data structure, and suddenly the thing that ``worked'' doesn't.
Lattice takes testing seriously: the language has built-in syntax for tests, a CLI runner that discovers and executes them, and a full suite of assertion functions that produce clear failure messages.
No external framework to install, no configuration file to wrangle.
Tests live right next to the code they verify.

%% ─────────────────────────────────────────────────────────────
\section{\texttt{test "name" \{ body \}} --- Inline Tests}\label{sec:inline-tests}
\index{test blocks!syntax}

Lattice's \lstinline|test| keyword lets you write tests directly in your source files.
The syntax is:

\begin{lstlisting}[language=Lattice, caption={Basic test block syntax}]
test "addition works" {
    assert(2 + 2 == 4, "basic arithmetic")
}
\end{lstlisting}

\begin{defnbox}{Test Block}
A \emph{test block} is a top-level declaration of the form \lstinline|test "name" \{ body \}|.
The name is a string literal that describes what the test verifies.
The body contains any number of statements---typically assertions---that run when the test suite is invoked via \lstinline|clat test|.
Test blocks are \emph{completely ignored} during normal program execution (\lstinline|clat run|).
\end{defnbox}

Let's write a more realistic example.
Suppose we have a temperature conversion module:

\begin{lstlisting}[language=Lattice, caption={Temperature conversion with inline tests}]
fn celsius_to_fahrenheit(c: Float) -> Float {
    return c * 9.0 / 5.0 + 32.0
}

fn fahrenheit_to_celsius(f: Float) -> Float {
    return (f - 32.0) * 5.0 / 9.0
}

test "freezing point conversion" {
    assert_eq(celsius_to_fahrenheit(0.0), 32.0)
    assert_eq(fahrenheit_to_celsius(32.0), 0.0)
}

test "boiling point conversion" {
    assert_eq(celsius_to_fahrenheit(100.0), 212.0)
    assert_eq(fahrenheit_to_celsius(212.0), 100.0)
}

test "body temperature" {
    let body_temp_f = celsius_to_fahrenheit(37.0)
    // 37.0 C should be approximately 98.6 F
    assert(body_temp_f > 98.5, "too low")
    assert(body_temp_f < 98.7, "too high")
}
\end{lstlisting}

The tests sit right below the functions they exercise.
Anyone reading this file immediately understands what the functions are expected to do.

\subsection{Tests Are Top-Level Declarations}\label{sec:tests-top-level}

Test blocks are parsed as top-level items by the Lattice parser (handled in \filename{src/parser.c} as \lstinline|ITEM_TEST|).
They cannot appear inside functions, loops, or other blocks.
This keeps them at the module level, visible and discoverable.

\begin{lstlisting}[language=Lattice, caption={Tests belong at the top level}]
fn compute(x: Int) -> Int {
    return x * x
}

// This is correct --- test at the top level:
test "compute squares" {
    assert_eq(compute(5), 25)
    assert_eq(compute(0), 0)
    assert_eq(compute(-3), 9)
}
\end{lstlisting}

\subsection{Tests Can Use Everything in the File}\label{sec:tests-scope}

When \lstinline|clat test| runs, it first executes all non-test top-level code (imports, function definitions, struct declarations, variable bindings) before running any test block.
This means tests have access to every function, struct, enum, and global defined in the file.

\begin{lstlisting}[language=Lattice, caption={Tests have access to all module-level definitions}]
struct Color {
    r: Int,
    g: Int,
    b: Int
}

fn mix(c1: Color, c2: Color) -> Color {
    return Color {
        r: (c1.r + c2.r) / 2,
        g: (c1.g + c2.g) / 2,
        b: (c1.b + c2.b) / 2
    }
}

test "mixing colors" {
    let red = Color { r: 255, g: 0, b: 0 }
    let blue = Color { r: 0, g: 0, b: 255 }
    let purple = mix(red, blue)

    assert_eq(purple.r, 127)
    assert_eq(purple.g, 0)
    assert_eq(purple.b, 127)
}
\end{lstlisting}

\begin{notebox}{Tests Run in Their Own Scope}
Each test block runs in a fresh scope.
Variables defined inside one test are not visible to other tests.
This prevents accidental coupling between test cases---each test is independent.
\end{notebox}

\subsection{What Happens During Normal Execution?}\label{sec:tests-normal-exec}

When you run a file with \lstinline|clat run|, test blocks are skipped entirely.
The compiler (in \filename{src/stackcompiler.c}) treats \lstinline|ITEM_TEST| items as no-ops during normal compilation.
This means you can sprinkle tests throughout your source code without any runtime cost in production.

\begin{lstlisting}[language=Lattice, caption={Tests are invisible during normal execution}]
fn greet(name: String) -> String {
    return "Hello, " + name + "!"
}

test "greeting includes the name" {
    assert_contains(greet("Lattice"), "Lattice")
}

// When running: clat run greeter.lat
// Only the code below executes:
print(greet("World"))
// Output: Hello, World!
\end{lstlisting}

%% ─────────────────────────────────────────────────────────────
\section{\texttt{clat test} --- Running the Test Suite}\label{sec:clat-test}
\index{clat test}

The \lstinline|clat test| command discovers and runs tests.
At its most basic:

\begin{lstlisting}[language=Lattice, caption={Running tests in a single file}]
// $ clat test temperature.lat
//   PASS  freezing point conversion
//   PASS  boiling point conversion
//   PASS  body temperature
//
// Results: 3 passed, 0 failed (3 total)
\end{lstlisting}

\subsection{Directory Discovery}\label{sec:test-discovery}
\index{clat test!discovery}

When you pass a directory instead of a file, \lstinline|clat test| recursively searches for test files.
The naming convention for test files is:

\begin{itemize}
  \item Files matching \lstinline|test_*.lat| (e.g., \lstinline|test_math.lat|, \lstinline|test_parser.lat|)
  \item Files matching \lstinline|*_test.lat| (e.g., \lstinline|math_test.lat|, \lstinline|parser_test.lat|)
\end{itemize}

\begin{lstlisting}[language=Lattice, caption={Running all tests in a directory}]
// $ clat test tests/
//   PASS  addition
//   PASS  subtraction
//   PASS  multiplication
//   FAIL  division by zero -- assertion failed: expected error
//
// Results: 3 passed, 1 failed (4 total)
\end{lstlisting}

The discovery is recursive, so tests nested in subdirectories are found too:

\begin{lstlisting}[language=Lattice, caption={Typical test directory layout}]
// tests/
//   test_math.lat
//   test_strings.lat
//   integration/
//     test_server.lat
//     test_database.lat
\end{lstlisting}

All four files would be discovered and executed.

\subsection{Filtering Tests by Name}\label{sec:test-filter}
\index{clat test!filter}

The \lstinline|-f| or \lstinline|--filter| flag runs only tests whose name contains the given substring:

\begin{lstlisting}[language=Lattice, caption={Filtering tests by name}]
// $ clat test tests/ -f "arithmetic"
//   PASS  basic arithmetic
//   PASS  floating-point arithmetic
//
// Results: 2 passed, 0 failed, 5 skipped (7 total)
\end{lstlisting}

The filter applies a simple substring match (using C's \lstinline|strstr|, as seen in the evaluator's \lstinline|evaluator_run_tests()| function in \filename{src/eval.c}).
If no tests match the filter, the runner reports how many were skipped.

\subsection{Verbose Mode}\label{sec:test-verbose}
\index{clat test!verbose}

Add \lstinline|-v| or \lstinline|--verbose| for more detailed output, including timing and setup information:

\begin{lstlisting}[language=Lattice, caption={Verbose test output}]
// $ clat test -v tests/test_math.lat
\end{lstlisting}

\subsection{Summary Mode}\label{sec:test-summary}

The \lstinline|--summary| flag produces a concise summary after the test run, particularly useful in CI environments where you want a quick overview.

\subsection{Full CLI Reference}\label{sec:test-cli-ref}

Here is the complete set of \lstinline|clat test| flags:

\begin{longtable}{@{}ll@{}}
\toprule
\textbf{Flag} & \textbf{Description} \\
\midrule
\endhead
\lstinline|-f, --filter <pattern>| & Only run tests whose name matches pattern \\
\lstinline|-v, --verbose| & Verbose output \\
\lstinline|--summary| & Show summary after run \\
\lstinline|--gc-stress| & GC stress testing mode \\
\lstinline|--no-regions| & Disable region-based memory \\
\lstinline|--no-assertions| & Disable runtime assertions \\
\lstinline|-h, --help| & Show help \\
\bottomrule
\end{longtable}

\begin{tipbox}{Exit Codes}
\lstinline|clat test| exits with code 0 if all tests pass, and code 1 if any test fails.
This makes it straightforward to integrate into CI pipelines and shell scripts.
\end{tipbox}

%% ─────────────────────────────────────────────────────────────
\section{Assert Functions}\label{sec:assert-functions}
\index{assert}\index{assert\_eq}\index{assert\_type}\index{assertions}

Lattice provides a rich set of built-in assertion functions.
These are not library functions---they're built into the evaluator and the runtime (\filename{src/eval.c}, \filename{src/runtime.c}), so they're always available without any imports.

\subsection{\texttt{assert(condition, message?)}}\label{sec:assert-basic}
\index{assert}

The fundamental assertion.
If the condition is falsy, the test fails with the optional message:

\begin{lstlisting}[language=Lattice, caption={Basic assert usage}]
test "assert basics" {
    assert(1 == 1, "one equals one")
    assert(true, "true is truthy")
    assert(42, "non-zero integers are truthy")
    assert("hello", "non-empty strings are truthy")
}
\end{lstlisting}

When \lstinline|assert| fails, it throws an error that the test runner catches.
The failure message includes your custom string, making it clear which assertion broke:

\begin{lstlisting}[language=Lattice, caption={A failing assert}]
test "this will fail" {
    let balance = -50
    assert(balance >= 0, "balance must not be negative")
    // Output: FAIL  this will fail -- assertion failed: balance must not be negative
}
\end{lstlisting}

\subsection{\texttt{assert\_eq(actual, expected)}}\label{sec:assert-eq}
\index{assert\_eq}

The workhorse of testing.
Checks that two values are equal, and if not, shows both values in the failure message:

\begin{lstlisting}[language=Lattice, caption={assert\_eq with clear failure messages}]
test "string operations" {
    let greeting = "hello" + " " + "world"
    assert_eq(greeting, "hello world")

    let trimmed = "  spaces  ".trim()
    assert_eq(trimmed, "spaces")
}

test "array operations" {
    let numbers = [1, 2, 3, 4, 5]
    assert_eq(len(numbers), 5)
    assert_eq(numbers[0], 1)
    assert_eq(numbers[4], 5)
}
\end{lstlisting}

If the assertion fails, you get a message like:
\begin{lstlisting}[language=Lattice]
// FAIL  string operations -- assert_eq failed: expected "spaces", got "  spaces  "
\end{lstlisting}

This level of detail in failure messages is worth its weight in gold during debugging.
You don't need to add print statements---the assertion tells you exactly what went wrong.

\subsection{\texttt{assert\_ne(actual, expected)}}\label{sec:assert-ne}
\index{assert\_ne}

The opposite of \lstinline|assert_eq|.
Fails if the two values are equal:

\begin{lstlisting}[language=Lattice, caption={Asserting inequality}]
test "unique ID generation" {
    let id1 = generate_id()
    let id2 = generate_id()
    assert_ne(id1, id2)
}
\end{lstlisting}

\subsection{\texttt{assert\_true(value)} and \texttt{assert\_false(value)}}\label{sec:assert-bool}
\index{assert\_true}\index{assert\_false}

Strict boolean checks.
Unlike \lstinline|assert|, which accepts any truthy value, \lstinline|assert_true| expects exactly \lstinline|true| and \lstinline|assert_false| expects exactly \lstinline|false|:

\begin{lstlisting}[language=Lattice, caption={Strict boolean assertions}]
test "boolean checks" {
    let is_admin = check_role("admin")
    assert_true(is_admin)

    let is_expired = check_expiry("2099-12-31")
    assert_false(is_expired)
}
\end{lstlisting}

\begin{warnbox}{assert\_true Is Strict}
\lstinline|assert_true(42)| will \emph{fail} because \lstinline|42| is not the boolean value \lstinline|true|, even though it's truthy.
If you want truthy checking, use \lstinline|assert(42, "expected truthy")|.
\end{warnbox}

\subsection{\texttt{assert\_type(value, type\_name)}}\label{sec:assert-type}
\index{assert\_type}

Verifies that a value has the expected type.
The type name is a string that matches the output of \lstinline|typeof()|, such as \lstinline|"Int"|, \lstinline|"Float"|, \lstinline|"String"|, \lstinline|"Array"|, or the name of a struct:

\begin{lstlisting}[language=Lattice, caption={Type assertions}]
struct Point {
    x: Float,
    y: Float
}

test "type checking" {
    assert_type(42, "Int")
    assert_type(3.14, "Float")
    assert_type("hello", "String")
    assert_type([1, 2, 3], "Array")
    assert_type(true, "Bool")

    let origin = Point { x: 0.0, y: 0.0 }
    assert_type(origin, "Point")
}
\end{lstlisting}

The \lstinline|assert_type| function checks both the primitive type and, for structs, the struct name.
So \lstinline|assert_type(origin, "Point")| succeeds because the evaluator (in \filename{src/eval.c}) compares the struct's name field against the expected string.

\subsection{\texttt{assert\_contains(haystack, needle)}}\label{sec:assert-contains}
\index{assert\_contains}

Checks that a string or array contains a given element:

\begin{lstlisting}[language=Lattice, caption={Containment assertions}]
test "string containment" {
    let message = "the quick brown fox"
    assert_contains(message, "brown")
    assert_contains(message, "fox")
}

test "array containment" {
    let colors = ["red", "green", "blue"]
    assert_contains(colors, "green")
}
\end{lstlisting}

\subsection{\texttt{assert\_throws(closure)}}\label{sec:assert-throws}
\index{assert\_throws}

Verifies that a piece of code throws an error.
Pass a closure that should fail:

\begin{lstlisting}[language=Lattice, caption={Asserting that code throws an error}]
test "division by zero throws" {
    assert_throws(|_| {
        let result = 10 / 0
    })
}

test "index out of bounds throws" {
    assert_throws(|_| {
        let items = [1, 2, 3]
        let bad = items[99]
    })
}
\end{lstlisting}

The built-in \lstinline|assert_throws| wraps the closure in a \lstinline|try|/\lstinline|catch| block.
If no error is thrown, the assertion fails with the message ``expected an error but none was thrown.''

\subsection{Summary of Built-in Assertions}\label{sec:assert-summary}

\begin{longtable}{@{}lp{8cm}@{}}
\toprule
\textbf{Function} & \textbf{Description} \\
\midrule
\endhead
\lstinline|assert(cond, msg?)| & Fails if \lstinline|cond| is falsy \\
\lstinline|assert_eq(a, b)| & Fails if \lstinline|a != b|, showing both values \\
\lstinline|assert_ne(a, b)| & Fails if \lstinline|a == b| \\
\lstinline|assert_true(v)| & Fails if \lstinline|v| is not exactly \lstinline|true| \\
\lstinline|assert_false(v)| & Fails if \lstinline|v| is not exactly \lstinline|false| \\
\lstinline|assert_type(v, t)| & Fails if \lstinline|typeof(v)| is not \lstinline|t| \\
\lstinline|assert_contains(h, n)| & Fails if string/array \lstinline|h| doesn't contain \lstinline|n| \\
\lstinline|assert_throws(fn)| & Fails if the closure does not throw \\
\bottomrule
\end{longtable}

%% ─────────────────────────────────────────────────────────────
\section{The Test Standard Library Module}\label{sec:test-stdlib}
\index{test standard library}\index{lib/test.lat}

Beyond the built-in assertions, Lattice ships with a test library at \filename{lib/test.lat} that provides a more structured, framework-style testing experience.
This is useful for larger projects that need test suites organized into groups.

\subsection{Importing the Test Library}\label{sec:test-import}

\begin{lstlisting}[language=Lattice, caption={Importing the test library}]
import "lib/test" as t
\end{lstlisting}

\subsection{Organizing Tests with \texttt{describe} and \texttt{it}}\label{sec:describe-it}
\index{describe}\index{it}

The test library provides two organizational primitives.
The \lstinline|describe()| function creates a named test suite, and \lstinline|it()| creates individual test cases within a suite.

\begin{lstlisting}[language=Lattice, caption={Structured test suites}]
import "lib/test" as t

t.run([
    t.describe("String operations", |_| {
        return [
            t.it("concatenation", |_| {
                t.assert_eq("hello" + " " + "world", "hello world")
            }),
            t.it("length", |_| {
                t.assert_eq(len("lattice"), 7)
            }),
            t.it("uppercase", |_| {
                t.assert_eq("hello".upper(), "HELLO")
            })
        ]
    }),
    t.describe("Math operations", |_| {
        return [
            t.it("addition", |_| {
                t.assert_eq(2 + 2, 4)
                t.assert_eq(1.5 + 2.5, 4.0)
            }),
            t.it("division by zero", |_| {
                t.assert_throws(|_| { 1 / 0 })
            })
        ]
    })
])
\end{lstlisting}

Running this with \lstinline|clat run| (not \lstinline|clat test|---this is a program, not a test file) produces:

\begin{lstlisting}[language=Lattice, caption={Structured test output}]
// Running tests...
//
//   String operations
//     checkmark concatenation
//     checkmark length
//     checkmark uppercase
//
//   Math operations
//     checkmark addition
//     checkmark division by zero
//
// 5 passed, 0 failed, 5 total
// Completed in 2ms
\end{lstlisting}

\begin{notebox}{Built-in Tests vs.\ the Test Library}
The built-in \lstinline|test "name" \{ ... \}| blocks are ideal for unit tests that live alongside source code, run by \lstinline|clat test|.
The \lstinline|lib/test| library is better for standalone test programs that need more structure---nested suites, timing, and a custom runner.
Both approaches work well; use whichever fits your project.
\end{notebox}

\subsection{Extended Assertions in the Test Library}\label{sec:test-lib-assertions}
\index{assert\_gt}\index{assert\_lt}\index{assert\_near}\index{assert\_nil}

The test library provides additional assertion functions beyond the built-ins:

\begin{lstlisting}[language=Lattice, caption={Extended assertions from lib/test}]
import "lib/test" as t

t.run([
    t.describe("Comparisons", |_| {
        return [
            t.it("greater than / less than", |_| {
                t.assert_gt(10, 5)
                t.assert_lt(3, 7)
                t.assert_gte(5, 5)
                t.assert_lte(5, 5)
            }),
            t.it("floating-point proximity", |_| {
                let pi = 3.14159265
                t.assert_near(pi, 3.14, 0.01)
            }),
            t.it("nil checks", |_| {
                t.assert_nil(nil)
                t.assert_not_nil(42)
                t.assert_not_nil("hello")
            }),
            t.it("type checking", |_| {
                t.assert_type(42, "Int")
                t.assert_type("hi", "String")
            })
        ]
    })
])
\end{lstlisting}

Here's the full list of assertions provided by \filename{lib/test.lat}:

\begin{longtable}{@{}lp{8cm}@{}}
\toprule
\textbf{Function} & \textbf{Description} \\
\midrule
\endhead
\lstinline|t.check(cond)| & Fail if condition is falsy \\
\lstinline|t.assert_eq(actual, expected)| & Fail if values differ \\
\lstinline|t.assert_neq(actual, expected)| & Fail if values are equal \\
\lstinline|t.assert_gt(a, b)| & Fail if \lstinline|a <= b| \\
\lstinline|t.assert_lt(a, b)| & Fail if \lstinline|a >= b| \\
\lstinline|t.assert_gte(a, b)| & Fail if \lstinline|a < b| \\
\lstinline|t.assert_lte(a, b)| & Fail if \lstinline|a > b| \\
\lstinline|t.assert_near(actual, expected, epsilon)| & Fail if values differ by more than \lstinline|epsilon| \\
\lstinline|t.assert_contains(haystack, needle)| & Fail if haystack doesn't contain needle \\
\lstinline|t.assert_throws(closure)| & Fail if the closure doesn't throw \\
\lstinline|t.assert_type(value, type_name)| & Fail if the type doesn't match \\
\lstinline|t.assert_nil(value)| & Fail if value is not \lstinline|nil| \\
\lstinline|t.assert_not_nil(value)| & Fail if value is \lstinline|nil| \\
\lstinline|t.assert_true(value)| & Fail if value is not \lstinline|true| \\
\lstinline|t.assert_false(value)| & Fail if value is not \lstinline|false| \\
\bottomrule
\end{longtable}

\begin{tipbox}{assert\_near for Floating-Point Tests}
Comparing floating-point numbers with \lstinline|assert_eq| can lead to spurious failures due to rounding.
Use \lstinline|t.assert_near(actual, expected, epsilon)| instead, where \lstinline|epsilon| is the acceptable margin of error.
For most calculations, an epsilon of \lstinline|0.0001| is a good starting point.
\end{tipbox}

\subsection{How the Test Library Works Internally}\label{sec:test-lib-internals}

If you open \filename{lib/test.lat}, you'll find that the library is written entirely in Lattice.
There's no magic---it uses the built-in \lstinline|assert()| function combined with \lstinline|format()| and \lstinline|repr()| to produce clear error messages, and \lstinline|try|/\lstinline|catch| to recover from assertion failures.

The \lstinline|it()| function returns a \lstinline|Map| with \lstinline|"name"| and \lstinline|"fn"| keys.
The \lstinline|describe()| function calls the builder closure to get an array of these Maps and wraps them in a suite Map.
The \lstinline|run()| function iterates through suites and test cases, invoking each closure inside a \lstinline|try|/\lstinline|catch| block.
If the closure completes without error, the test passes.
If an assertion triggers an error, the catch block captures it and reports the failure.

%% ─────────────────────────────────────────────────────────────
\section{Testing Strategies for Lattice Programs}\label{sec:testing-strategies}
\index{testing strategies}

Having the tools is one thing.
Using them effectively is another.
Let's explore strategies that experienced Lattice developers rely on.

\subsection{Unit Testing: One Function, One Test}\label{sec:unit-testing}

The most common pattern: write a test block for each public function.
Keep tests close to the code they exercise.

\begin{lstlisting}[language=Lattice, caption={Unit tests alongside the code}]
fn is_palindrome(s: String) -> Bool {
    let reversed = s.chars().reverse().join("")
    return s == reversed
}

test "is_palindrome with palindromes" {
    assert_true(is_palindrome("racecar"))
    assert_true(is_palindrome("level"))
    assert_true(is_palindrome("a"))
    assert_true(is_palindrome(""))
}

test "is_palindrome with non-palindromes" {
    assert_false(is_palindrome("hello"))
    assert_false(is_palindrome("lattice"))
}
\end{lstlisting}

\subsection{Testing Edge Cases}\label{sec:edge-cases}

The most productive tests often check boundaries: empty inputs, zero values, negative numbers, maximum sizes, and \lstinline|nil|.

\begin{lstlisting}[language=Lattice, caption={Testing edge cases}]
fn safe_divide(a: Float, b: Float) -> Float {
    if b == 0.0 {
        return 0.0
    }
    return a / b
}

test "safe_divide normal cases" {
    assert_eq(safe_divide(10.0, 2.0), 5.0)
    assert_eq(safe_divide(7.0, 3.5), 2.0)
}

test "safe_divide edge cases" {
    assert_eq(safe_divide(0.0, 5.0), 0.0)
    assert_eq(safe_divide(10.0, 0.0), 0.0)
    assert_eq(safe_divide(0.0, 0.0), 0.0)
    assert_eq(safe_divide(-10.0, 2.0), -5.0)
}
\end{lstlisting}

\subsection{Testing Error Paths}\label{sec:testing-errors}

Use \lstinline|assert_throws| (either the built-in or the library version) to verify that functions fail when they should:

\begin{lstlisting}[language=Lattice, caption={Testing that errors are thrown correctly}]
fn parse_positive_int(s: String) -> Int {
    let n = parse_int(s)
    if is_error(n) {
        return error("not a valid integer: " + s)
    }
    if n <= 0 {
        return error("expected positive integer, got " + to_string(n))
    }
    return n
}

test "parse_positive_int rejects invalid input" {
    assert_throws(|_| { parse_positive_int("abc") })
    assert_throws(|_| { parse_positive_int("-5") })
    assert_throws(|_| { parse_positive_int("0") })
}

test "parse_positive_int accepts valid input" {
    assert_eq(parse_positive_int("42"), 42)
    assert_eq(parse_positive_int("1"), 1)
}
\end{lstlisting}

\subsection{Testing Structs and Methods}\label{sec:testing-structs}

When testing structs, create instances with known values and verify field access, method results, and trait implementations:

\begin{lstlisting}[language=Lattice, caption={Testing structs}]
struct Rectangle {
    width: Float,
    height: Float
}

fn area(r: Rectangle) -> Float {
    return r.width * r.height
}

fn perimeter(r: Rectangle) -> Float {
    return 2.0 * (r.width + r.height)
}

fn is_square(r: Rectangle) -> Bool {
    return r.width == r.height
}

test "rectangle area" {
    let r = Rectangle { width: 5.0, height: 3.0 }
    assert_eq(area(r), 15.0)
}

test "rectangle perimeter" {
    let r = Rectangle { width: 5.0, height: 3.0 }
    assert_eq(perimeter(r), 16.0)
}

test "square detection" {
    let square = Rectangle { width: 4.0, height: 4.0 }
    let rect = Rectangle { width: 4.0, height: 6.0 }
    assert_true(is_square(square))
    assert_false(is_square(rect))
}
\end{lstlisting}

\subsection{Organizing a Test Suite}\label{sec:organizing-tests}

For a project of any size, keep tests in a dedicated \filename{tests/} directory.
Each test file mirrors the module it tests:

\begin{lstlisting}[language=Lattice, caption={Project layout with tests}]
// src/
//   math.lat
//   strings.lat
//   http.lat
// tests/
//   test_math.lat
//   test_strings.lat
//   test_http.lat
\end{lstlisting}

Then run the entire suite with:

\begin{lstlisting}[language=Lattice, caption={Running the full test suite}]
// $ clat test tests/
\end{lstlisting}

\subsection{The Test-First Workflow}\label{sec:test-first}

A powerful practice: write the test \emph{before} the implementation.

\begin{lstlisting}[language=Lattice, caption={Writing the test first}]
// Step 1: Write the test
fn fibonacci(n: Int) -> Int {
    // TODO: implement
    return 0
}

test "fibonacci sequence" {
    assert_eq(fibonacci(0), 0)
    assert_eq(fibonacci(1), 1)
    assert_eq(fibonacci(2), 1)
    assert_eq(fibonacci(5), 5)
    assert_eq(fibonacci(10), 55)
}

// Step 2: Run it (it fails)
// $ clat test fib.lat
//   FAIL  fibonacci sequence -- assert_eq failed: expected 1, got 0

// Step 3: Implement until it passes
\end{lstlisting}

The failing test gives you a clear target.
Once all assertions pass, you know the implementation is correct---at least for the cases you specified.

\subsection{Testing with GC Stress Mode}\label{sec:gc-stress}
\index{gc-stress}\index{testing!GC stress}

The \lstinline|--gc-stress| flag triggers garbage collection after every allocation.
This is invaluable for catching memory bugs:

\begin{lstlisting}[language=Lattice, caption={Running tests with GC stress}]
// $ clat test --gc-stress tests/
\end{lstlisting}

Under GC stress, any test that accidentally holds a dangling reference or forgets to root a value will fail quickly.
It makes tests slower, so you wouldn't run this in every CI build---but running it periodically catches subtle bugs that normal execution misses.

\subsection{Disabling Assertions for Performance Testing}\label{sec:no-assertions}

The \lstinline|--no-assertions| flag disables runtime assertions (specifically, the \lstinline|debug_assert| built-in).
This is useful when you want to benchmark code that includes debug assertions without removing them from the source:

\begin{lstlisting}[language=Lattice, caption={Running without runtime assertions}]
// $ clat test --no-assertions benchmarks/
\end{lstlisting}

Note that \lstinline|assert()| (the testing assertion) is not affected by this flag---only \lstinline|debug_assert()| is suppressed.

%% ─────────────────────────────────────────────────────────────
\section*{Exercises}\label{sec:testing-exercises}
\addcontentsline{toc}{section}{Exercises}

\begin{enumerate}
  \item Write a \lstinline|stack.lat| module that implements a stack using an array (with \lstinline|push|, \lstinline|pop|, \lstinline|peek|, and \lstinline|is_empty| functions). Include inline \lstinline|test| blocks that verify each operation, including edge cases like popping from an empty stack.

  \item Create a structured test suite using \lstinline|lib/test| that tests a string utility library with at least three \lstinline|describe| groups and five total test cases. Use \lstinline|assert_eq|, \lstinline|assert_contains|, and \lstinline|assert_throws|.

  \item Write a test file that uses \lstinline|assert_near| to verify floating-point calculations (trigonometry, compound interest, or physics formulas). Experiment with different epsilon values to understand how precision affects test results.

  \item Set up a small project with three source files in \filename{src/} and corresponding test files in \filename{tests/}. Run the entire test suite, then use \lstinline|--filter| to run only the tests related to one module. Write a one-line CI command that formats and tests the whole project.

  \item Write a test that verifies the \emph{behavior} of the phase system. For example, verify that freezing a value with \lstinline|freeze| prevents mutation by catching the resulting error with \lstinline|assert_throws|.
\end{enumerate}

%% ─────────────────────────────────────────────────────────────
\section*{What's Next}\label{sec:testing-next}
\addcontentsline{toc}{section}{What's Next}

When a test fails and the assertion message isn't enough to pinpoint the problem, you need a more powerful tool: a debugger.
In \cref{ch:debugging}, we'll explore Lattice's built-in CLI debugger, learn how to set breakpoints, step through code line by line, inspect variables, and even connect VS Code for a full graphical debugging experience.
